\section{O-ISAC FOUNDATIONS AND PERFORMANCE COMPARISON}
\label{sec:foundations_comparison}
A useful comparison of O-ISAC designs asks how a shared design choice affects
both data transmission and sensing. Answering this question requires an
understanding of the signal paths, the quantities being measured, and the
conditions under which the results were obtained. This section explains these
connections and establishes how we compare performance throughout the review.

\begin{figure*}[!t]
\centering
\footnotesize
\begin{tabular}{@{}c@{\hspace{9pt}}c@{\hspace{9pt}}c@{}}
\fbox{\parbox[c][13mm][c]{0.28\textwidth}{\centering
\textbf{Shared design factor}\\
A variable, constraint, or disturbance}}
& $\longrightarrow$ &
\begin{tabular}{@{}c@{}}
\fbox{\parbox[c][9mm][c]{0.43\textwidth}{\centering
\textbf{Communication result}\\
Rate, error rate, or another communication measure}}\\[5pt]
\fbox{\parbox[c][9mm][c]{0.43\textwidth}{\centering
\textbf{Sensing result}\\
Error, detection probability, or another sensing measure}}
\end{tabular}
\end{tabular}
\caption{Connecting communication and sensing performance within a study.
Both results are read against the same design factor, with their measurement
conditions and configuration identified. The source must show how both
results relate to that factor. Comparison with another design requires the
additional checks in Fig.~\ref{fig:claim_permission_pipeline}.}
\label{fig:joint_evidence_anatomy}
\end{figure*}

\begin{figure*}[!t]
\centering
\footnotesize
\begin{tabular}{@{}c@{\hspace{7pt}}c@{\hspace{7pt}}c@{\hspace{7pt}}c@{\hspace{7pt}}c@{}}
\fbox{\parbox[c][14mm][c]{0.255\textwidth}{\centering
\textbf{Do the metrics mean the same thing?}\\
Definition, measurement point, and units}}
& $\longrightarrow$ &
\fbox{\parbox[c][14mm][c]{0.255\textwidth}{\centering
\textbf{Are the settings comparable?}\\
Physical system, task, conditions, and baseline}}
& $\longrightarrow$ &
\fbox{\parbox[c][14mm][c]{0.255\textwidth}{\centering
\textbf{Does the evidence support the comparison?}\\
Independent studies, validation, and uncertainty}}
\end{tabular}
\caption{Checks for a numerical comparison across studies. All three must
support the comparison being made. Table~\ref{tab:comparison_record} summarizes
the conclusions that different kinds of evidence can support.}
\label{fig:claim_permission_pipeline}
\end{figure*}

\begin{table*}[!t]
\caption{What different kinds of O-ISAC evidence can tell us}
\label{tab:comparison_record}
\centering
\footnotesize
\setlength{\tabcolsep}{4pt}
\renewcommand{\arraystretch}{1.15}
\begin{tabularx}{\textwidth}{@{}>{\raggedright\arraybackslash}p{0.18\textwidth}
>{\raggedright\arraybackslash}p{0.38\textwidth}
>{\raggedright\arraybackslash}X@{}}
\toprule
Evidence & What it establishes & How we use it \\
\midrule
Shared architecture &
Which components, signals, or resources serve both functions. &
Explain the integration mechanism. Joint performance requires results
that connect a shared factor to both functions. \\
\addlinespace[3pt]
An individual result &
Communication or sensing performance under identified conditions. &
Report what the system achieved for that function and setting. \\
\addlinespace[3pt]
Linked results within one study &
How both functions relate to a shared factor under the same or matched
conditions, supported by a common configuration or timing record. &
Explain the performance relationship and any tradeoff demonstrated within
the tested range. \\
\addlinespace[3pt]
Comparable results from independent studies &
A numerical comparison supported by the metric, setting, and evidence
checks in Fig.~\ref{fig:claim_permission_pipeline}. &
Compare the reported performance within the conditions that the studies
have in common. \\
\addlinespace[3pt]
A recurring mechanism &
A relationship supported in several study settings whose numerical results
cannot be directly combined. &
Explain the mechanism and any supported direction of change.
Keep numerical magnitudes tied to their source settings. \\
\addlinespace[3pt]
A recurring obstacle to comparison &
Which missing measurement, baseline, or validation prevents a useful
engineering conclusion. &
Specify an experiment or reporting requirement that could resolve the
question. \\
\bottomrule
\end{tabularx}
\end{table*}

\subsection{Linking Communication and Sensing Performance}

The physical signal path determines how communication and sensing interact.
A fiber can serve as the sensing medium or carry a signal to another segment
where sensing takes place. In photonic-terahertz systems, optical components
may generate or process the carrier, while a radio-frequency path interacts
with the target. We therefore distinguish the sensing modality from the
hardware used to produce or process the signal
\cite{OISAC_SCR00220,OISAC_SCR00940,OISAC_SCR00196,OISAC_SCR00083,
OISAC_SCR00366}. Section~\ref{sec:optical_modality_families} develops these
physical differences for each platform.

Within a system, communication and sensing may share hardware, a signal path,
or a waveform. They can also be linked through resource allocation or a common
design objective. Several forms of sharing can coexist. Identifying what is
shared explains the architecture. Evaluating its performance requires a
further question about how a design variable, constraint, or disturbance is
linked to both functions.

For example, to examine the effect of power allocation, we need communication
and sensing results for corresponding allocation settings. Those results can
show whether a gain in one function is accompanied by an improvement, a cost,
or little measured change in the other. The relationship must be traceable
to the same study configuration or to settings whose differences can be
accounted for from the source. Two performance values reported separately in a
paper do not, by themselves, establish that relationship.

Figure~\ref{fig:joint_evidence_anatomy} illustrates the connection we seek in
each study. A relationship established in one study supports conclusions
within its reported conditions. Applying it to another design also requires
an understanding of the physical paths and shared resources in both systems,
followed by the measurement and comparison checks below.

\subsection{Understanding the Reported Measurements}

A metric becomes useful for comparison when its meaning and measurement
point are clear. Optical signal-to-noise ratio (OSNR) and electrical
signal-to-noise ratio (SNR) refer to optical and electrical measurement
domains, respectively.
Launch power and received power likewise refer to different locations along
the link. These distinctions determine which quantities can be compared
directly and which require a source-supported conversion.

Communication metrics also depend on how performance is counted. Gross rate
and net rate use different accounting conventions, as do error rates before
and after forward error correction. Sensing metrics distinguish different
abilities. Resolution concerns the separation of targets, while estimation
error describes deviation from a reference value. Precision, detection
probability, and analytical bounds provide further information that must be
interpreted according to the source definition.

Operating conditions complete the interpretation. We retain the task and
target together with the channel, geometry, and processing settings.
The baseline identifies what a reported improvement is measured against, and
the ground truth provides the reference for sensing accuracy. Communication
and sensing results are linked only when their settings are the same or can
be matched using information reported in the source. If those conditions
differ materially, we discuss each
result within its own setting.

The form of validation also matters. An analytical bound, a simulation, and a
measurement on deployed hardware answer different questions about performance.
For each result, we retain its validation setting and exact source location.
To connect the two functions, we also follow the model configuration,
experimental run, or recorded timing through which their results were obtained.
Field deployment alone does not establish that both functions were evaluated
together. The conclusion follows the conditions actually supported by the
study.

\subsection{Comparing Results and Drawing Design Lessons}

A within-study comparison shows how communication and sensing respond to a
shared design factor over the tested operating range. We report the direction
of that relationship and any magnitude established by the source. Its
relevance to another design depends on the physical context, the shared
mechanism, and the measurement conditions.

Numerical comparisons across studies require at least two independent studies
with compatible physical systems and tasks. For a numerical comparison of
joint performance, each study must connect the shared factor to both
communication and sensing results. We check that metric definitions,
measurement points, and units agree. We then examine operating conditions and
baselines, followed by validation methods and relevant uncertainty.
Figure~\ref{fig:claim_permission_pipeline} groups these checks.
Any remaining differences must be small enough, or sufficiently accounted
for, to support the particular comparison.

When numerical values cannot be compared on this basis, studies may still
reveal a recurring design mechanism. We can explain how that mechanism is
linked to the two functions while keeping each numerical result in its own
setting. We do not combine these values into a common effect or use them to
rank platforms. Any conflicting evidence that changes the interpretation is
discussed alongside the conclusion.

Table~\ref{tab:comparison_record} summarizes how we use the evidence. A useful
research recommendation also follows this reasoning. It identifies a
recurring obstacle to comparison and specifies the measurement, reference
experiment, or validation needed to address it. This makes the recommendation
a testable next step.

Section~\ref{sec:review_methods} explains how the study records support this
analysis. Section~\ref{sec:optical_modality_families} develops the architectures,
and Section~\ref{sec:tradeoffs} examines performance and tradeoffs.
Sections~\ref{sec:validation_reproducibility}
and~\ref{sec:technologies_applications_6g} assess validation and application
requirements. These findings guide the research priorities in
Section~\ref{sec:discussion_roadmap}.

\FloatBarrier
